\section{Conclusion}
\label{sec:conclusion}
With leakage kept out of every fitted step, a single re-centred transfer configuration let stroke
patients start quickly on a sixth-finger BCI: after one trial of each movement, three-class
accuracy was \rqTwoRestAhSfRctTslrKOneMean\,\% with transfer against
\rqTwoRestAhSfTsLrKOneMean\,\% without it, after which the transfer curve barely moved, and with
the whole pool (\rqTwoMinutesKAll\,min nominal, \rqTwoMinutesRecordedKAll\,min as recorded) the
decoder trained on the patient alone came out numerically, though not significantly, ahead. The
data did not show a sixth-finger command of its own. What they support offline is imagined
movement against rest, a possible ``move now'' switch, which in a continuous replay caught
\exFiveFingerFingerOnMean\,\% of sixth-finger periods with \exFiveFingerFalsePerMinMean{} false
activations a minute. The ablations, mostly exploratory, suggest that re-centring at small
budgets, the choice of channels and the continuous gate shaped these results most, while crops,
mirroring and healthy grasp imagery did not help. Testing the switch in closed loop is the next step.
